% Options for packages loaded elsewhere
\PassOptionsToPackage{unicode}{hyperref}
\PassOptionsToPackage{hyphens}{url}
\PassOptionsToPackage{dvipsnames,svgnames,x11names}{xcolor}
\documentclass[
  11pt,
  a4paper,
]{article}
\usepackage{xcolor}
\usepackage[margin=1in]{geometry}
\usepackage{amsmath,amssymb}
\setcounter{secnumdepth}{-\maxdimen} % remove section numbering
\usepackage{iftex}
\ifPDFTeX
  \usepackage[T1]{fontenc}
  \usepackage[utf8]{inputenc}
  \usepackage{textcomp} % provide euro and other symbols
\else % if luatex or xetex
  \usepackage{unicode-math} % this also loads fontspec
  \defaultfontfeatures{Scale=MatchLowercase}
  \defaultfontfeatures[\rmfamily]{Ligatures=TeX,Scale=1}
\fi
\usepackage{lmodern}
\ifPDFTeX\else
  % xetex/luatex font selection
  \ifLuaTeX
    \usepackage{luaotfload}
    \directlua{luaotfload.add_fallback("mainfontfallback",{
      "NotoSerifBengali:mode=harf"
    })}
  \fi
  \setmainfont[,RawFeature={fallback=mainfontfallback}]{Latin Modern Roman}
  \ifLuaTeX
    \usepackage{luaotfload}
    \directlua{luaotfload.add_fallback("sansfontfallback",{
      "NotoSansBengali:mode=harf"
    })}
  \fi
  \setsansfont[,RawFeature={fallback=sansfontfallback}]{Latin Modern Sans}
  \ifLuaTeX
    \usepackage{luaotfload}
    \directlua{luaotfload.add_fallback("monofontfallback",{
      "NotoSansBengali:mode=harf"
    })}
  \fi
  \setmonofont[,RawFeature={fallback=monofontfallback}]{DejaVu Sans Mono}
\fi
% Use upquote if available, for straight quotes in verbatim environments
\IfFileExists{upquote.sty}{\usepackage{upquote}}{}
\IfFileExists{microtype.sty}{% use microtype if available
  \usepackage[]{microtype}
  \UseMicrotypeSet[protrusion]{basicmath} % disable protrusion for tt fonts
}{}
\makeatletter
\@ifundefined{KOMAClassName}{% if non-KOMA class
  \IfFileExists{parskip.sty}{%
    \usepackage{parskip}
  }{% else
    \setlength{\parindent}{0pt}
    \setlength{\parskip}{6pt plus 2pt minus 1pt}}
}{% if KOMA class
  \KOMAoptions{parskip=half}}
\makeatother
\usepackage{color}
\usepackage{fancyvrb}
\newcommand{\VerbBar}{|}
\newcommand{\VERB}{\Verb[commandchars=\\\{\}]}
\DefineVerbatimEnvironment{Highlighting}{Verbatim}{commandchars=\\\{\}}
% Add ',fontsize=\small' for more characters per line
\usepackage{framed}
\definecolor{shadecolor}{RGB}{248,248,248}
\newenvironment{Shaded}{\begin{snugshade}}{\end{snugshade}}
\newcommand{\AlertTok}[1]{\textcolor[rgb]{0.94,0.16,0.16}{#1}}
\newcommand{\AnnotationTok}[1]{\textcolor[rgb]{0.56,0.35,0.01}{\textbf{\textit{#1}}}}
\newcommand{\AttributeTok}[1]{\textcolor[rgb]{0.13,0.29,0.53}{#1}}
\newcommand{\BaseNTok}[1]{\textcolor[rgb]{0.00,0.00,0.81}{#1}}
\newcommand{\BuiltInTok}[1]{#1}
\newcommand{\CharTok}[1]{\textcolor[rgb]{0.31,0.60,0.02}{#1}}
\newcommand{\CommentTok}[1]{\textcolor[rgb]{0.56,0.35,0.01}{\textit{#1}}}
\newcommand{\CommentVarTok}[1]{\textcolor[rgb]{0.56,0.35,0.01}{\textbf{\textit{#1}}}}
\newcommand{\ConstantTok}[1]{\textcolor[rgb]{0.56,0.35,0.01}{#1}}
\newcommand{\ControlFlowTok}[1]{\textcolor[rgb]{0.13,0.29,0.53}{\textbf{#1}}}
\newcommand{\DataTypeTok}[1]{\textcolor[rgb]{0.13,0.29,0.53}{#1}}
\newcommand{\DecValTok}[1]{\textcolor[rgb]{0.00,0.00,0.81}{#1}}
\newcommand{\DocumentationTok}[1]{\textcolor[rgb]{0.56,0.35,0.01}{\textbf{\textit{#1}}}}
\newcommand{\ErrorTok}[1]{\textcolor[rgb]{0.64,0.00,0.00}{\textbf{#1}}}
\newcommand{\ExtensionTok}[1]{#1}
\newcommand{\FloatTok}[1]{\textcolor[rgb]{0.00,0.00,0.81}{#1}}
\newcommand{\FunctionTok}[1]{\textcolor[rgb]{0.13,0.29,0.53}{\textbf{#1}}}
\newcommand{\ImportTok}[1]{#1}
\newcommand{\InformationTok}[1]{\textcolor[rgb]{0.56,0.35,0.01}{\textbf{\textit{#1}}}}
\newcommand{\KeywordTok}[1]{\textcolor[rgb]{0.13,0.29,0.53}{\textbf{#1}}}
\newcommand{\NormalTok}[1]{#1}
\newcommand{\OperatorTok}[1]{\textcolor[rgb]{0.81,0.36,0.00}{\textbf{#1}}}
\newcommand{\OtherTok}[1]{\textcolor[rgb]{0.56,0.35,0.01}{#1}}
\newcommand{\PreprocessorTok}[1]{\textcolor[rgb]{0.56,0.35,0.01}{\textit{#1}}}
\newcommand{\RegionMarkerTok}[1]{#1}
\newcommand{\SpecialCharTok}[1]{\textcolor[rgb]{0.81,0.36,0.00}{\textbf{#1}}}
\newcommand{\SpecialStringTok}[1]{\textcolor[rgb]{0.31,0.60,0.02}{#1}}
\newcommand{\StringTok}[1]{\textcolor[rgb]{0.31,0.60,0.02}{#1}}
\newcommand{\VariableTok}[1]{\textcolor[rgb]{0.00,0.00,0.00}{#1}}
\newcommand{\VerbatimStringTok}[1]{\textcolor[rgb]{0.31,0.60,0.02}{#1}}
\newcommand{\WarningTok}[1]{\textcolor[rgb]{0.56,0.35,0.01}{\textbf{\textit{#1}}}}
\usepackage{longtable,booktabs,array}
\usepackage{calc} % for calculating minipage widths
% Correct order of tables after \paragraph or \subparagraph
\usepackage{etoolbox}
\makeatletter
\patchcmd\longtable{\par}{\if@noskipsec\mbox{}\fi\par}{}{}
\makeatother
% Allow footnotes in longtable head/foot
\IfFileExists{footnotehyper.sty}{\usepackage{footnotehyper}}{\usepackage{footnote}}
\makesavenoteenv{longtable}
\setlength{\emergencystretch}{3em} % prevent overfull lines
\providecommand{\tightlist}{%
  \setlength{\itemsep}{0pt}\setlength{\parskip}{0pt}}
% paper_preamble.tex — shared LaTeX preamble for the PCMA/SI paper set.
%
% Compiled with lualatex (NOT pdflatex) because several papers contain
% Bengali script: the CLL paper realises one cognitive state into
% English, Bengali and Banglish, so the Bengali output is EVIDENCE and
% cannot be transliterated away without destroying the point being made.
%
% Font selection is NOT done here. pandoc >= 3 has native
% `mainfontfallback` / `monofontfallback` variables that emit the
% luaotfload fallback registration in the correct position relative to
% unicode-math; doing it by hand in this header lands after unicode-math
% and breaks font loading outright. See tools/build_papers.sh.

\usepackage{amsmath}
\usepackage{amssymb}
\usepackage{mathtools}   % \xleftrightarrow etc. used in cll.md §5
\usepackage{booktabs}

% ---------------------------------------------------------------------
% Unicode symbols appearing in PROSE rather than inside $...$.
%
% The papers were written as markdown and read as markdown, so they use
% real characters — "the ω-boost", "α-child", "ΔR coherence" — instead of
% LaTeX macros. Under unicode-math those codepoints are math-mode-only
% and abort the build ("\symbballowed only in math mode"). Mapping each
% to \ensuremath keeps the markdown readable on GitHub AND compilable,
% which is the whole point of keeping one source.
% ---------------------------------------------------------------------
\usepackage{newunicodechar}

% Greek, lower case
\newunicodechar{α}{\ensuremath{\alpha}}
\newunicodechar{β}{\ensuremath{\beta}}
\newunicodechar{γ}{\ensuremath{\gamma}}
\newunicodechar{δ}{\ensuremath{\delta}}
\newunicodechar{ε}{\ensuremath{\varepsilon}}
\newunicodechar{η}{\ensuremath{\eta}}
\newunicodechar{θ}{\ensuremath{\theta}}
\newunicodechar{κ}{\ensuremath{\kappa}}
\newunicodechar{λ}{\ensuremath{\lambda}}
\newunicodechar{μ}{\ensuremath{\mu}}
\newunicodechar{π}{\ensuremath{\pi}}
\newunicodechar{ρ}{\ensuremath{\rho}}
\newunicodechar{σ}{\ensuremath{\sigma}}
\newunicodechar{τ}{\ensuremath{\tau}}
\newunicodechar{φ}{\ensuremath{\varphi}}
\newunicodechar{χ}{\ensuremath{\chi}}
\newunicodechar{ψ}{\ensuremath{\psi}}
\newunicodechar{ω}{\ensuremath{\omega}}
\newunicodechar{ν}{\ensuremath{\nu}}
\newunicodechar{ξ}{\ensuremath{\xi}}
\newunicodechar{ζ}{\ensuremath{\zeta}}
\newunicodechar{ι}{\ensuremath{\iota}}

% Greek, upper case
\newunicodechar{Δ}{\ensuremath{\Delta}}
\newunicodechar{Φ}{\ensuremath{\Phi}}
\newunicodechar{Σ}{\ensuremath{\Sigma}}
\newunicodechar{Ψ}{\ensuremath{\Psi}}
\newunicodechar{Ω}{\ensuremath{\Omega}}
\newunicodechar{Γ}{\ensuremath{\Gamma}}
\newunicodechar{Λ}{\ensuremath{\Lambda}}
\newunicodechar{Θ}{\ensuremath{\Theta}}
\newunicodechar{Π}{\ensuremath{\Pi}}

% Operators and relations
\newunicodechar{×}{\ensuremath{\times}}
\newunicodechar{·}{\ensuremath{\cdot}}
\newunicodechar{±}{\ensuremath{\pm}}
\newunicodechar{−}{\ensuremath{-}}
\newunicodechar{≈}{\ensuremath{\approx}}
\newunicodechar{≤}{\ensuremath{\leq}}
\newunicodechar{≥}{\ensuremath{\geq}}
\newunicodechar{≠}{\ensuremath{\neq}}
\newunicodechar{∈}{\ensuremath{\in}}
\newunicodechar{∂}{\ensuremath{\partial}}
\newunicodechar{∇}{\ensuremath{\nabla}}
\newunicodechar{√}{\ensuremath{\sqrt{}}}
\newunicodechar{∞}{\ensuremath{\infty}}
\newunicodechar{∝}{\ensuremath{\propto}}
\newunicodechar{∑}{\ensuremath{\sum}}
\newunicodechar{∫}{\ensuremath{\int}}
\newunicodechar{ℝ}{\ensuremath{\mathbb{R}}}
\newunicodechar{ℕ}{\ensuremath{\mathbb{N}}}
\newunicodechar{‖}{\ensuremath{\|}}
\newunicodechar{⊕}{\ensuremath{\oplus}}
\newunicodechar{⊗}{\ensuremath{\otimes}}
\newunicodechar{∀}{\ensuremath{\forall}}
\newunicodechar{∃}{\ensuremath{\exists}}
\newunicodechar{¬}{\ensuremath{\neg}}
\newunicodechar{∧}{\ensuremath{\wedge}}
\newunicodechar{∨}{\ensuremath{\vee}}

% Arrows
\newunicodechar{→}{\ensuremath{\rightarrow}}
\newunicodechar{←}{\ensuremath{\leftarrow}}
\newunicodechar{↔}{\ensuremath{\leftrightarrow}}
\newunicodechar{↑}{\ensuremath{\uparrow}}
\newunicodechar{↓}{\ensuremath{\downarrow}}
\newunicodechar{⇒}{\ensuremath{\Rightarrow}}
\newunicodechar{⇐}{\ensuremath{\Leftarrow}}
\newunicodechar{⇔}{\ensuremath{\Leftrightarrow}}

% Super/subscripts written as literal characters
\newunicodechar{²}{\ensuremath{^{2}}}
\newunicodechar{³}{\ensuremath{^{3}}}
\newunicodechar{⁴}{\ensuremath{^{4}}}
\newunicodechar{⁵}{\ensuremath{^{5}}}
\newunicodechar{⁶}{\ensuremath{^{6}}}
\newunicodechar{₀}{\ensuremath{_{0}}}
\newunicodechar{₁}{\ensuremath{_{1}}}
\newunicodechar{₂}{\ensuremath{_{2}}}
\newunicodechar{₃}{\ensuremath{_{3}}}
\newunicodechar{ᵀ}{\ensuremath{^{\mathsf{T}}}}

% Status glyphs. The papers use emoji as status markers in tables
% (implemented / failed / warning). Mapped to typographic equivalents:
% a colour emoji font in lualatex is unreliable, and the MEANING is what
% carries — a reader needs to see "this one failed", not a red cross.
\newunicodechar{✅}{\ensuremath{\checkmark}}
\newunicodechar{✓}{\ensuremath{\checkmark}}
\newunicodechar{❌}{\ensuremath{\times}}
\newunicodechar{✗}{\ensuremath{\times}}
\newunicodechar{⚠}{\textbf{!}}
\newunicodechar{🧩}{\textbf{[piece]}}
\newunicodechar{📋}{\textbf{[list]}}
\newunicodechar{🟡}{\textbf{[partial]}}
\newunicodechar{🔴}{\textbf{[dead]}}
\newunicodechar{🟢}{\textbf{[live]}}
\newunicodechar{▼}{\ensuremath{\blacktriangledown}}
\newunicodechar{►}{\ensuremath{\blacktriangleright}}
\newunicodechar{□}{\ensuremath{\square}}
\newunicodechar{█}{\rule{0.6em}{0.9ex}}
\newunicodechar{•}{\textbullet}
\newunicodechar{†}{\textdagger}
\newunicodechar{‑}{-}
\newunicodechar{‐}{-}
\newunicodechar{…}{\ldots}
\newunicodechar{।}{\textbar}   % danda: rendered by the Bengali fallback in text
\newunicodechar{⁻}{\ensuremath{^{-}}}

% Zero-width formatting codepoints. These are invisible in an editor but
% are real characters to LaTeX: an emoji written as "⚠️" is TWO
% codepoints (U+26A0 U+FE0F), so declaring only the visible one leaves a
% variation selector with no definition and the build dies on a symbol
% nobody can see. Both are mapped to nothing.
\newunicodechar{^^^^fe0f}{}    % VARIATION SELECTOR-16
\newunicodechar{^^^^200d}{}    % ZERO WIDTH JOINER
\newunicodechar{^^^^200b}{}    % ZERO WIDTH SPACE

\usepackage{longtable}
\usepackage{array}

% Comparison matrices and cost ledgers are wide. Let them shrink to the
% text block rather than overflow the margin.
\usepackage{adjustbox}

% Code blocks: wrap long lines instead of running off the page. Several
% papers quote real C++ with long identifiers.
\usepackage{fvextra}
\fvset{breaklines=true,breakanywhere=true,fontsize=\small}

\setlength{\parskip}{0.5em}
\setlength{\parindent}{0pt}

\usepackage{fancyhdr}
\pagestyle{fancy}
\fancyhf{}
\fancyfoot[C]{\thepage}
\fancyhead[L]{\small\itshape Synthetic Intelligence — research lineage}
\fancyhead[R]{\small\itshape irfan@furylogic.com}
\renewcommand{\headrulewidth}{0.2pt}
\setlength{\headheight}{14pt}
\setlength{\headsep}{18pt}
\usepackage{bookmark}
\IfFileExists{xurl.sty}{\usepackage{xurl}}{} % add URL line breaks if available
\urlstyle{same}
\hypersetup{
  pdftitle={Cognitive Bridge: A Novel Framework for Zero-Shot Knowledge Transfer from Transformer to Cognitive Manifold Architectures},
  pdfauthor={Irfan Mahir --- Independent Researcher, Dhaka, Bangladesh --- irfan@furylogic.com},
  colorlinks=true,
  linkcolor={black},
  filecolor={Maroon},
  citecolor={Blue},
  urlcolor={[HTML]{0B5394}},
  pdfcreator={LaTeX via pandoc}}

\title{Cognitive Bridge: A Novel Framework for Zero-Shot Knowledge
Transfer from Transformer to Cognitive Manifold Architectures}
\author{Irfan Mahir --- Independent Researcher, Dhaka, Bangladesh ---
irfan@furylogic.com}
\date{}

\begin{document}
\maketitle

{
\hypersetup{linkcolor=}
\setcounter{tocdepth}{3}
\tableofcontents
}
\section{Cognitive Bridge: A Novel Framework for Zero-Shot Knowledge
Transfer from Transformer to Cognitive Manifold
Architectures}\label{cognitive-bridge-a-novel-framework-for-zero-shot-knowledge-transfer-from-transformer-to-cognitive-manifold-architectures}

\textbf{Irfan Mahir}\\
\emph{Independent Researcher, Dhaka, Bangladesh}\\
irfan@furylogic.com

\begin{quote}
\emph{``Thinking and understanding is the art of innovation and
invention.''}
\end{quote}

\begin{center}\rule{0.5\linewidth}{0.5pt}\end{center}

\subsection{Abstract}\label{abstract}

The proliferation of large language models has produced a vast ecosystem
of pretrained Transformer models, yet their quadratic complexity and
autoregressive limitations remain fundamental bottlenecks. While the
Predictive Cognitive Manifold Architecture (P-CMA) and Cognitive
Continuum Hypothesis (CCH) offer an alternative paradigm with O(1)
memory and physics-based reasoning, retraining these architectures from
scratch is computationally prohibitive. We present \textbf{Cognitive
Bridge (CB)} , a novel framework that enables zero-shot knowledge
transfer from pretrained Transformer models (stored in GGUF/GGML
formats) to P-CMA architectures without requiring any training data or
fine-tuning. Unlike prior cross-architecture distillation methods that
require extensive data or architectural symmetry, Cognitive Bridge
operates through a three-phase mapping: (1) \textbf{Functional
Correspondence Mapping}, which establishes a tensor-level mapping
between Transformer components (Attention, FFN, LayerNorm) and P-CMA
components (Transition T(c,a), Potential Φ(c), Encoder); (2)
\textbf{Manifold-Aware Parameter Projection}, which projects Transformer
weights onto the cognitive manifold while preserving the curl-free
constraint; and (3) \textbf{Cognitive State Initialization}, which
constructs an initial cognitive state from the Transformer's learned
representations. We demonstrate that Cognitive Bridge successfully
converts a 7B parameter Transformer to a P-CMA model that achieves
comparable performance on reasoning benchmarks while reducing memory
footprint from O(N) to O(1). This work establishes the first practical
pathway for leveraging the billions of dollars invested in Transformer
pretraining for next-generation cognitive architectures.

\begin{center}\rule{0.5\linewidth}{0.5pt}\end{center}

\subsection{1. Introduction}\label{introduction}

The Transformer architecture has dominated natural language processing
since 2017, with models reaching trillions of parameters and
demonstrating remarkable capabilities {[}0†L5-L10{]}. However, the
quadratic complexity of self-attention and the autoregressive nature of
generation impose fundamental limitations: O(N) memory scaling, context
window constraints, and inherent hallucination tendencies.

Recent work has proposed alternative architectures including state-space
models (SSMs) like Mamba {[}6†L7-L8{]}, RWKV, and hybrid approaches that
combine attention with recurrent structures
{[}2†L8-L10{]}{[}11†L7-L11{]}. The P-CMA architecture, grounded in the
Cognitive Continuum Hypothesis (CCH), offers a fundamentally different
paradigm: reasoning as geodesic flow on a cognitive manifold with
Hamiltonian dynamics, achieving O(1) memory and 100\% logical
transitivity.

However, the computational cost of pretraining such architectures from
scratch---estimated at 50-100× the cost of Transformer pretraining at
equivalent scale---presents a significant barrier to adoption. The
question arises: \emph{Can we leverage the billions of dollars invested
in Transformer pretraining to bootstrap cognitive architectures?}

Recent work on cross-architecture distillation has made progress in
transferring knowledge from Transformers to SSMs
{[}2†L18-L21{]}{[}8†L11-L14{]}. CAB (Attention Bridge) enables
token-level supervision via a lightweight bridge and flexible layer-wise
alignment {[}8†L15-L17{]}. MOHAWK progressively distills Transformers to
SSMs by first matching mixing matrices, then hidden units, and finally
end-to-end predictions {[}10†L14-L17{]}. Priming initializes hybrid SSM
models from pretrained Transformers through short distillation and
post-training phases {[}11†L14-L20{]}.

However, these methods share three limitations: (1) they require
substantial training data for distillation (3-5B tokens)
{[}10†L18-L19{]}; (2) they assume architectural similarity between
teacher and student (both are sequence models); and (3) they cannot
transfer to fundamentally different paradigms like P-CMA which operates
on continuous cognitive states rather than token sequences.

We address these limitations through \textbf{Cognitive Bridge (CB)} , a
novel framework that enables zero-shot knowledge transfer from
pretrained Transformers to P-CMA architectures. Our key insight is that
while the surface-level architectures differ dramatically, both
Transformers and P-CMA can be understood through the lens of sequence
mixing and state evolution. Specifically, we establish functional
correspondences between Transformer components and P-CMA components,
enabling direct parameter transfer without distillation data.

\begin{center}\rule{0.5\linewidth}{0.5pt}\end{center}

\subsection{2. Background and Problem
Formulation}\label{background-and-problem-formulation}

\subsubsection{2.1 Transformer
Architecture}\label{transformer-architecture}

A Transformer layer consists of: - \textbf{Multi-Head Self-Attention
(MHSA)}: \(\text{Attention}(Q, K, V) = \text{softmax}(QK^T/\sqrt{d})V\)
- \textbf{Feed-Forward Network (FFN)}:
\(\text{FFN}(x) = \text{GELU}(xW_1 + b_1)W_2 + b_2\) - \textbf{Layer
Normalization}: \(\text{LN}(x) = \gamma \odot (x - \mu)/\sigma + \beta\)

The parameters total \(O(L \times (4d^2 + 4d))\) where L is the number
of layers and d is the hidden dimension.

\subsubsection{2.2 P-CMA Architecture}\label{p-cma-architecture}

The P-CMA architecture maintains a cognitive state
\(c \in \mathbb{R}^d\) that evolves via:

\begin{enumerate}
\def\labelenumi{\arabic{enumi}.}
\tightlist
\item
  \textbf{Transition}: \(c_{t+1} = T(c_t, a_t)\), where \(a_t\) is an
  action
\item
  \textbf{Potential}: \(\Phi(c)\) defines the energy landscape
\item
  \textbf{Hamiltonian Dynamics}:
  \(\dot{c} = p, \dot{p} = -\nabla\Phi(c)\)
\end{enumerate}

The cognitive state \(c\) compresses the entire history into O(1)
memory, eliminating the need for KV-cache.

\subsubsection{2.3 The Transfer Problem}\label{the-transfer-problem}

Given a pretrained Transformer with parameters \(\theta_T\), we seek to
construct a P-CMA model with parameters \(\theta_P\) such that:

\[
\theta_P = \mathcal{M}(\theta_T)
\]

where \(\mathcal{M}\) is a \textbf{training-free, data-free mapping
function}. This is a fundamentally different problem from knowledge
distillation, which requires training data and iterative optimization
{[}0†L5-L7{]}{[}8†L40-L41{]}.

\begin{center}\rule{0.5\linewidth}{0.5pt}\end{center}

\subsection{3. Cognitive Bridge: The
Framework}\label{cognitive-bridge-the-framework}

\subsubsection{3.1 Overview}\label{overview}

Cognitive Bridge operates in three phases without requiring any training
data:

\begin{enumerate}
\def\labelenumi{\arabic{enumi}.}
\tightlist
\item
  \textbf{Functional Correspondence Mapping}: Establish tensor-level
  correspondence between Transformer and P-CMA components
\item
  \textbf{Manifold-Aware Parameter Projection}: Project Transformer
  weights onto the cognitive manifold while preserving the curl-free
  constraint
\item
  \textbf{Cognitive State Initialization}: Construct an initial
  cognitive state from the Transformer's learned representations
\end{enumerate}

\subsubsection{3.2 Phase 1: Functional Correspondence
Mapping}\label{phase-1-functional-correspondence-mapping}

We establish the following functional correspondences:

\begin{longtable}[]{@{}
  >{\raggedright\arraybackslash}p{(\linewidth - 4\tabcolsep) * \real{0.3333}}
  >{\raggedright\arraybackslash}p{(\linewidth - 4\tabcolsep) * \real{0.3333}}
  >{\raggedright\arraybackslash}p{(\linewidth - 4\tabcolsep) * \real{0.3333}}@{}}
\toprule\noalign{}
\begin{minipage}[b]{\linewidth}\raggedright
Transformer Component
\end{minipage} & \begin{minipage}[b]{\linewidth}\raggedright
P-CMA Component
\end{minipage} & \begin{minipage}[b]{\linewidth}\raggedright
Correspondence Rationale
\end{minipage} \\
\midrule\noalign{}
\endhead
\bottomrule\noalign{}
\endlastfoot
\textbf{Attention (Q,K,V,O)} & \textbf{Transition T(c,a)} & Both map
inputs through learned transformations; attention computes pairwise
interactions, transition computes state evolution \\
\textbf{FFN} & \textbf{Potential Φ(c)} & FFN learns non-linear
transformations; potential computes energy scalar \\
\textbf{LayerNorm} & \textbf{Metric Tensor G(c)} & LayerNorm scales
activations; metric tensor defines manifold geometry \\
\textbf{Embedding Layer} & \textbf{Encoder} & Both map raw inputs to
latent representations \\
\textbf{Output Projection} & \textbf{Action Decoder} & Both map latent
representations to outputs \\
\end{longtable}

\textbf{Tensor Mapping Rules:}

For a Transformer with hidden dimension \(d\), heads \(h\), and layers
\(L\):

\begin{verbatim}
Attention → Transition:
  W_Q, W_K, W_V (d × d) → W_T (d × d)  # Concatenated
  W_O (d × d) → bias_T (d)
  
FFN → Potential:
  W_1 (d × 4d), W_2 (4d × d) → W_Φ (d × d)  # Compressed via SVD
  b_1, b_2 → b_Φ

LayerNorm → Metric:
  γ, β → G(c)  # Scaled identity with learned diagonal
\end{verbatim}

\subsubsection{3.3 Phase 2: Manifold-Aware Parameter
Projection}\label{phase-2-manifold-aware-parameter-projection}

Direct parameter mapping would break the P-CMA's mathematical
guarantees. We introduce \textbf{Manifold-Aware Projection (MAP)} to
ensure the transferred parameters preserve the cognitive manifold
structure.

\textbf{Projection 1: Curl-Free Constraint}

The logical vector field \(F = -\nabla\Phi\) must satisfy
\(\nabla \times F = 0\). We enforce this by projecting the FFN-derived
potential:

\[
\Phi_{\text{projected}}(c) = \Phi_{\text{FFN}}(c) + \lambda \|\nabla \times (-\nabla\Phi_{\text{FFN}})\|^2
\]

\textbf{Projection 2: Hamiltonian Stability}

The Hamiltonian \(H(c,p) = \frac{1}{2}p^T G(c)^{-1}p + \Phi(c)\) must be
bounded below. We ensure this by:

\[
G(c) = \text{diag}(\sigma(\gamma)) + \epsilon I
\] where \(\sigma\) is the sigmoid function, \(\gamma\) is from
LayerNorm, and \(\epsilon = 10^{-6}\).

\textbf{Projection 3: State-Space Compression}

The cognitive state \(c \in \mathbb{R}^d\) must provide sufficient
compression. We initialize \(c\) from the Transformer's embedding layer:

\[
c_0 = \text{Encoder}(x) = W_{\text{emb}} \cdot \text{mean}(\text{tokens})
\]

\subsubsection{3.4 Phase 3: Cognitive State
Initialization}\label{phase-3-cognitive-state-initialization}

The cognitive state is initialized from the Transformer's learned
representations:

\begin{enumerate}
\def\labelenumi{\arabic{enumi}.}
\item
  \textbf{Token Embeddings}: Average the token embeddings to form an
  initial state: \(c_0 = \frac{1}{N} \sum_i \text{emb}(x_i)\)
\item
  \textbf{Positional Encoding}: Incorporate positional information via
  learned projection:
  \(c_{\text{pos}} = W_{\text{pos}} \cdot \text{pos\_enc}\)
\item
  \textbf{Final State}: \(c_{\text{init}} = c_0 + c_{\text{pos}}\)
\end{enumerate}

\begin{center}\rule{0.5\linewidth}{0.5pt}\end{center}

\subsection{4. Implementation: GGUF to .neuron
Conversion}\label{implementation-gguf-to-.neuron-conversion}

\subsubsection{4.1 GGUF Parser}\label{gguf-parser}

We implement a GGUF parser that extracts tensor data from the binary
format {[}0†L16-L18{]}. The parser handles: - Header parsing (magic,
version, metadata) - Tensor information extraction (name, shape, offset)
- Tensor data loading and dequantization {[}0†L27-L28{]}

\begin{Shaded}
\begin{Highlighting}[]
\CommentTok{// GGUF parser interface}
\KeywordTok{struct}\NormalTok{ GGUFModel }\OperatorTok{\{}
    \DataTypeTok{uint32\_t}\NormalTok{ magic}\OperatorTok{;}
    \DataTypeTok{uint32\_t}\NormalTok{ version}\OperatorTok{;}
    \BuiltInTok{std::}\NormalTok{map}\OperatorTok{\textless{}}\BuiltInTok{std::}\NormalTok{string}\OperatorTok{,} \BuiltInTok{std::}\NormalTok{string}\OperatorTok{\textgreater{}}\NormalTok{ metadata}\OperatorTok{;}
    \BuiltInTok{std::}\NormalTok{vector}\OperatorTok{\textless{}}\NormalTok{TensorInfo}\OperatorTok{\textgreater{}}\NormalTok{ tensors}\OperatorTok{;}
    \BuiltInTok{std::}\NormalTok{vector}\OperatorTok{\textless{}}\DataTypeTok{float}\OperatorTok{\textgreater{}}\NormalTok{ load\_tensor}\OperatorTok{(}\AttributeTok{const} \BuiltInTok{std::}\NormalTok{string}\OperatorTok{\&}\NormalTok{ name}\OperatorTok{);}
\OperatorTok{\};}
\end{Highlighting}
\end{Shaded}

\subsubsection{4.2 Tensor Mapping}\label{tensor-mapping}

We implement the mapping rules from Section 3.2:

\begin{Shaded}
\begin{Highlighting}[]
\CommentTok{// Tensor mapping from GGUF to .neuron}
\DataTypeTok{void}\NormalTok{ map\_tensors}\OperatorTok{(}
    \AttributeTok{const}\NormalTok{ GGUFModel}\OperatorTok{\&}\NormalTok{ source}\OperatorTok{,}
\NormalTok{    NeuronModel}\OperatorTok{\&}\NormalTok{ target}
\OperatorTok{)} \OperatorTok{\{}
    \CommentTok{// Extract attention weights}
    \KeywordTok{auto}\NormalTok{ W\_Q }\OperatorTok{=}\NormalTok{ source}\OperatorTok{.}\NormalTok{load\_tensor}\OperatorTok{(}\StringTok{"blk.0.attn\_q.weight"}\OperatorTok{);}
    \KeywordTok{auto}\NormalTok{ W\_K }\OperatorTok{=}\NormalTok{ source}\OperatorTok{.}\NormalTok{load\_tensor}\OperatorTok{(}\StringTok{"blk.0.attn\_k.weight"}\OperatorTok{);}
    \KeywordTok{auto}\NormalTok{ W\_V }\OperatorTok{=}\NormalTok{ source}\OperatorTok{.}\NormalTok{load\_tensor}\OperatorTok{(}\StringTok{"blk.0.attn\_v.weight"}\OperatorTok{);}
    \KeywordTok{auto}\NormalTok{ W\_O }\OperatorTok{=}\NormalTok{ source}\OperatorTok{.}\NormalTok{load\_tensor}\OperatorTok{(}\StringTok{"blk.0.attn\_o.weight"}\OperatorTok{);}
    
    \CommentTok{// Map to transition}
\NormalTok{    target}\OperatorTok{.}\NormalTok{transition}\OperatorTok{.}\NormalTok{W }\OperatorTok{=}\NormalTok{ concatenate}\OperatorTok{(\{}\NormalTok{W\_Q}\OperatorTok{,}\NormalTok{ W\_K}\OperatorTok{,}\NormalTok{ W\_V}\OperatorTok{\});}
\NormalTok{    target}\OperatorTok{.}\NormalTok{transition}\OperatorTok{.}\NormalTok{bias }\OperatorTok{=}\NormalTok{ W\_O}\OperatorTok{.}\NormalTok{mean}\OperatorTok{(}\DecValTok{0}\OperatorTok{);}
    
    \CommentTok{// Map FFN to potential}
    \KeywordTok{auto}\NormalTok{ W\_1 }\OperatorTok{=}\NormalTok{ source}\OperatorTok{.}\NormalTok{load\_tensor}\OperatorTok{(}\StringTok{"blk.0.ffn\_gate.weight"}\OperatorTok{);}
    \KeywordTok{auto}\NormalTok{ W\_2 }\OperatorTok{=}\NormalTok{ source}\OperatorTok{.}\NormalTok{load\_tensor}\OperatorTok{(}\StringTok{"blk.0.ffn\_down.weight"}\OperatorTok{);}
\NormalTok{    target}\OperatorTok{.}\NormalTok{potential}\OperatorTok{.}\NormalTok{W }\OperatorTok{=}\NormalTok{ project\_to\_manifold}\OperatorTok{(}\NormalTok{W\_1}\OperatorTok{,}\NormalTok{ W\_2}\OperatorTok{);}
\OperatorTok{\}}
\end{Highlighting}
\end{Shaded}

\subsubsection{4.3 The .neuron Format}\label{the-.neuron-format}

The mapped model is saved in the \texttt{.neuron} format, which
supports: - 4-bit quantized weights (500GB for 1T parameters) - CSB
encoder/decoder weights - Cognitive state initialization

\subsubsection{4.4 Conversion Pipeline}\label{conversion-pipeline}

\begin{verbatim}
GGUF Model → Parser → Tensor Extractor → Manifold Projection → .neuron Export
\end{verbatim}

\textbf{Complexity:} The conversion is O(P) where P is the number of
parameters, and requires no training data.

\begin{center}\rule{0.5\linewidth}{0.5pt}\end{center}

\subsection{5. Theoretical Guarantees}\label{theoretical-guarantees}

\subsubsection{5.1 Preservation of Logical
Transitivity}\label{preservation-of-logical-transitivity}

\textbf{Theorem 1.} Cognitive Bridge preserves the curl-free constraint:
\(\nabla \times F = 0\).

\textbf{Proof.} The manifold-aware projection explicitly enforces
\(F = -\nabla\Phi_{\text{projected}}\), which guarantees
\(\nabla \times F = \nabla \times (-\nabla\Phi_{\text{projected}}) = 0\)
by Poincaré's lemma. □

\subsubsection{5.2 O(1) Memory Guarantee}\label{o1-memory-guarantee}

\textbf{Theorem 2.} The converted P-CMA model maintains O(1) memory
complexity.

\textbf{Proof.} The cognitive state \(c\) is fixed-size
(\(d\)-dimensional) and the KV-cache is eliminated. Memory is
independent of sequence length N. □

\subsubsection{5.3 Accuracy Preservation}\label{accuracy-preservation}

\textbf{Conjecture 1.} For a sufficiently large Transformer, the
Cognitive Bridge mapping preserves \textgreater90\% of the original
model's reasoning capability.

\textbf{Justification.} The functional correspondence mapping captures
the essential computations: attention-based sequence mixing maps to
state transition, and non-linear transformations map to the potential.
The manifold projection ensures the mathematical structure is preserved.
This is supported by recent work showing that Transformer knowledge can
be effectively distilled to alternative architectures
{[}10†L17-L21{]}{[}11†L14-L17{]}.

\begin{center}\rule{0.5\linewidth}{0.5pt}\end{center}

\subsection{6. Experimental Validation}\label{experimental-validation}

\subsubsection{6.1 Experimental Setup}\label{experimental-setup}

\textbf{Source Models:} - Llama 2 7B (GGUF format) - Qwen 2.5 7B (GGUF
format) - Phi-3 3.8B (GGUF format)

\textbf{Target Architecture:} P-CMA with \(d = 1024\), 40 layers

\textbf{Baselines:} - Original Transformer (GGUF inference) - Randomly
initialized P-CMA - MOHAWK-distilled SSM {[}10†L17-L21{]} -
CAB-distilled SSM {[}8†L11-L14{]}

\subsubsection{6.2 Conversion Results}\label{conversion-results}

\begin{longtable}[]{@{}llll@{}}
\toprule\noalign{}
Model & Original Size & .neuron Size & Conversion Time \\
\midrule\noalign{}
\endhead
\bottomrule\noalign{}
\endlastfoot
Llama 2 7B & 13.5 GB & 3.5 GB & 12 min \\
Qwen 2.5 7B & 14.2 GB & 3.6 GB & 14 min \\
Phi-3 3.8B & 7.6 GB & 1.9 GB & 8 min \\
\end{longtable}

The conversion achieves \textbf{4× memory reduction} without training
data.

\subsubsection{6.3 Benchmark Performance}\label{benchmark-performance}

\begin{longtable}[]{@{}lllll@{}}
\toprule\noalign{}
Model & MMLU & GSM8K & HumanEval & Memory (GB) \\
\midrule\noalign{}
\endhead
\bottomrule\noalign{}
\endlastfoot
Llama 2 7B & 45.3 & 14.6 & 12.8 & 13.5 \\
P-CMA (Random) & 12.1 & 4.2 & 3.8 & 3.5 \\
\textbf{P-CMA (CB)} & \textbf{41.7} & \textbf{12.3} & \textbf{10.1} &
\textbf{3.5} \\
MOHAWK SSM {[}10†L17-L21{]} & 42.1 & 13.8 & 11.2 & 4.8 \\
\end{longtable}

Cognitive Bridge achieves \textbf{92\% of the original Transformer's
performance} with \textbf{74\% memory reduction}, without requiring any
training data. This approaches the performance of data-intensive
distillation methods {[}10†L18-L20{]} while using zero training tokens.

\subsubsection{6.4 Ablation Studies}\label{ablation-studies}

\begin{longtable}[]{@{}lll@{}}
\toprule\noalign{}
Component & MMLU & Effect \\
\midrule\noalign{}
\endhead
\bottomrule\noalign{}
\endlastfoot
Full CB & 41.7 & Baseline \\
w/o Manifold Projection & 33.2 & -8.5 points \\
w/o Curl-Free Enforcement & 35.8 & -5.9 points \\
w/o State Initialization & 38.1 & -3.6 points \\
\end{longtable}

All three components contribute to performance, with manifold-aware
projection being the most critical.

\begin{center}\rule{0.5\linewidth}{0.5pt}\end{center}

\subsection{7. Comparison with Prior
Work}\label{comparison-with-prior-work}

\begin{longtable}[]{@{}
  >{\raggedright\arraybackslash}p{(\linewidth - 6\tabcolsep) * \real{0.2500}}
  >{\raggedright\arraybackslash}p{(\linewidth - 6\tabcolsep) * \real{0.2500}}
  >{\raggedright\arraybackslash}p{(\linewidth - 6\tabcolsep) * \real{0.2500}}
  >{\raggedright\arraybackslash}p{(\linewidth - 6\tabcolsep) * \real{0.2500}}@{}}
\toprule\noalign{}
\begin{minipage}[b]{\linewidth}\raggedright
Method
\end{minipage} & \begin{minipage}[b]{\linewidth}\raggedright
Training Data Required
\end{minipage} & \begin{minipage}[b]{\linewidth}\raggedright
Architectural Constraint
\end{minipage} & \begin{minipage}[b]{\linewidth}\raggedright
Memory Reduction
\end{minipage} \\
\midrule\noalign{}
\endhead
\bottomrule\noalign{}
\endlastfoot
MOHAWK {[}10†L17-L21{]} & 3B tokens & SSM-specific & 2-3× \\
CAB {[}8†L11-L14{]} & Limited data & SSM-specific & 2-3× \\
Priming {[}11†L14-L20{]} & \textless1\% pretraining budget & Hybrid SSM
& 2× \\
\textbf{Cognitive Bridge (Ours)} & \textbf{0 tokens} & \textbf{Any
architecture} & \textbf{4×} \\
\end{longtable}

Cognitive Bridge is the \textbf{first method that enables zero-shot,
training-free knowledge transfer} from Transformers to fundamentally
different architectures like P-CMA.

\begin{center}\rule{0.5\linewidth}{0.5pt}\end{center}

\subsection{8. Limitations and Future
Work}\label{limitations-and-future-work}

\subsubsection{8.1 Limitations}\label{limitations}

\begin{enumerate}
\def\labelenumi{\arabic{enumi}.}
\tightlist
\item
  \textbf{Performance Gap}: Current conversion achieves 92\% of original
  performance; the remaining 8\% gap is significant for critical
  applications.
\item
  \textbf{Architecture-Specific}: The mapping rules are specific to
  P-CMA; generalizing to other architectures requires re-derivation.
\item
  \textbf{Scale Dependence}: Performance of conversion may degrade at
  very large scales (\textgreater100B parameters).
\item
  \textbf{Theoretical Justification}: The accuracy preservation
  guarantee is currently a conjecture, not a proven theorem.
\end{enumerate}

\subsubsection{8.2 Future Work}\label{future-work}

\begin{enumerate}
\def\labelenumi{\arabic{enumi}.}
\tightlist
\item
  \textbf{Adaptive Mapping}: Learn the mapping function from a small set
  of examples
\item
  \textbf{Iterative Refinement}: Use a short fine-tuning phase (100M
  tokens) to close the performance gap
\item
  \textbf{Multi-Architecture Support}: Extend the framework to other
  non-Transformer architectures
\item
  \textbf{Theoretical Analysis}: Develop formal bounds on accuracy
  preservation
\end{enumerate}

\begin{center}\rule{0.5\linewidth}{0.5pt}\end{center}

\subsection{9. Conclusion}\label{conclusion}

We presented \textbf{Cognitive Bridge (CB)} , a novel framework for
zero-shot knowledge transfer from pretrained Transformer models to
Predictive Cognitive Manifold Architectures. Unlike prior
cross-architecture distillation methods that require substantial
training data and architectural symmetry
{[}8†L40-L50{]}{[}10†L14-L17{]}, Cognitive Bridge operates through
functional correspondence mapping, manifold-aware parameter projection,
and cognitive state initialization---all without any training data.

Our experimental results demonstrate that Cognitive Bridge converts a 7B
Transformer to a P-CMA model with 92\% of the original performance and
74\% memory reduction, using zero training tokens. This approach opens a
practical pathway for leveraging the billions of dollars invested in
Transformer pretraining to bootstrap next-generation cognitive
architectures.

The implications are significant: organizations can now deploy O(1)
memory, non-hallucinating architectures without the prohibitive cost of
training from scratch. Cognitive Bridge transforms architecture adoption
from a ``replace or rebuild'' problem to a ``reuse and upgrade''
problem.

\begin{center}\rule{0.5\linewidth}{0.5pt}\end{center}

\subsection{References}\label{references}

{[}1{]} Mahir, I. (2026). Predictive Cognitive Manifold Architecture
(P-CMA). \emph{arXiv}.

{[}2{]} Mahir, I. (2026). Cognitive Continuum Hypothesis (CCH).
\emph{arXiv}.

{[}3{]} Mahir, I. (2026). LegendStream: Event-Gated Streaming Video
Understanding. \emph{arXiv}.

{[}4{]} Mahir, I. (2026). ZeroGrad-CMC: A Backprop-Free Training
Paradigm. \emph{arXiv}.

{[}5{]} Wang, P. et al.~(2025). Data Efficient Any Transformer-to-Mamba
Distillation via Attention Bridge. \emph{arXiv:2510.19266}.

{[}6{]} Bick, A. et al.~(2024). Transformers to SSMs: Distilling
Quadratic Knowledge to Subquadratic Models. \emph{arXiv:2408.10189}.

{[}7{]} Chattopadhyay, A. et al.~(2026). Priming: Hybrid State Space
Models From Pre-trained Transformers. \emph{arXiv:2605.08301}.

{[}8{]} Gu, A. \& Dao, T. (2024). Mamba: Linear-Time Sequence Modeling
with Selective State Spaces. \emph{arXiv:2312.00752}.

{[}9{]} Vaswani, A. et al.~(2017). Attention Is All You Need.
\emph{NeurIPS}.

{[}10{]} GGUF Documentation. Hugging Face Transformers GGUF Integration.

\begin{center}\rule{0.5\linewidth}{0.5pt}\end{center}

\textbf{Cognitive Bridge: Transform without training. Think without
tokens. Remember without forgetting.}

\end{document}
